\documentclass[12pt,a4paper]{article}

% ── packages ────────────────────────────────────────────────────────────────
\usepackage[utf8]{inputenc}
\usepackage[T1]{fontenc}
\usepackage{geometry}
\usepackage{graphicx}
\usepackage{booktabs}
\usepackage{array}
\usepackage{tabularx}
\usepackage{longtable}
\usepackage{xcolor}
\usepackage{colortbl}
\usepackage{fancyhdr}
\usepackage{titlesec}
\usepackage{parskip}
\usepackage{setspace}
\usepackage{caption}
\usepackage{float}
\usepackage{listings}
\usepackage{hyperref}

\geometry{
  top=2.5cm, bottom=2.5cm,
  left=3cm,  right=2cm
}

% ── colours ─────────────────────────────────────────────────────────────────
\definecolor{navy}{HTML}{0C447C}
\definecolor{headbg}{HTML}{185FA5}
\definecolor{rowalt}{HTML}{F8F8F6}
\definecolor{border}{HTML}{D3D1C7}
\definecolor{graydk}{HTML}{5F5E5A}
\definecolor{amberdk}{HTML}{633806}
\definecolor{purpledk}{HTML}{3C3489}
\definecolor{greendk}{HTML}{27500A}
\definecolor{sqlbg}{HTML}{F4F4F4}
\definecolor{sqlkw}{HTML}{0C447C}
\definecolor{sqlstr}{HTML}{3B6D11}

% ── hyperref ────────────────────────────────────────────────────────────────
\hypersetup{
  colorlinks=true,
  linkcolor=navy,
  urlcolor=navy,
  pdftitle={Modelagem SBD OLTP - Locadora de Veiculos},
  pdfauthor={Izabela Lima da Silva | Caio Meirelles}
}

% ── section formatting ───────────────────────────────────────────────────────
\titleformat{\section}
  {\large\bfseries\color{navy}}
  {\thesection.}{0.5em}{}[\vspace{-0.3em}\textcolor{border}{\rule{\linewidth}{0.4pt}}]

\titleformat{\subsection}
  {\normalsize\bfseries\color{navy}}
  {\thesubsection.}{0.5em}{}

% ── header / footer ─────────────────────────────────────────────────────────
\pagestyle{fancy}
\fancyhf{}
\renewcommand{\headrulewidth}{0pt}
\setlength{\headheight}{15pt}
\fancyhead[L]{\small\color{graydk} MAE016 TED-B \textbar\ Modelagem SBD OLTP}
\fancyhead[R]{\small\color{graydk} UFRJ 2026.1}
\fancyfoot[C]{\small\color{graydk}\thepage}

% ── table helpers ────────────────────────────────────────────────────────────
\newcolumntype{L}[1]{>{\raggedright\arraybackslash}p{#1}}
\newcolumntype{C}[1]{>{\centering\arraybackslash}p{#1}}

\newcommand{\badgePK}{\textcolor{headbg}{\textbf{PK}}}
\newcommand{\badgeFK}{\textcolor{amberdk}{\textbf{FK}}}
\newcommand{\badgeUQ}{\textcolor{purpledk}{\textbf{UQ}}}
\newcommand{\badgeNN}{\textcolor{greendk}{\textbf{NN}}}

% ── SQL listing style ────────────────────────────────────────────────────────
\lstdefinestyle{sql}{
  language=SQL,
  backgroundcolor=\color{sqlbg},
  basicstyle=\ttfamily\small,
  keywordstyle=\bfseries\color{sqlkw},
  stringstyle=\color{sqlstr},
  commentstyle=\color{graydk},
  frame=single,
  framerule=0.4pt,
  rulecolor=\color{border},
  breaklines=true,
  showstringspaces=false,
  tabsize=4,
  xleftmargin=0.5cm,
  xrightmargin=0.5cm,
  aboveskip=0.5em,
  belowskip=0.5em,
  morekeywords={VARCHAR,DATE,DECIMAL,REFERENCES,INT}
}

% ── document ─────────────────────────────────────────────────────────────────
\begin{document}

% ════════════════════════════════════════════════════════════════════════════
%  CAPA
% ════════════════════════════════════════════════════════════════════════════
\begin{titlepage}
  \thispagestyle{empty}
  \begin{center}

    {\color{navy}\rule{\linewidth}{2pt}}
    \vspace{0.4cm}

    {\Large\bfseries\color{navy}
      Universidade Federal do Rio de Janeiro\\[0.2em]
      N\'ucleo de Computa\c{c}\~ao Eletr\^onica --- NCE}

    \vspace{0.6cm}
    {\color{navy}\rule{\linewidth}{0.5pt}}
    \vspace{1.5cm}

    {\LARGE\bfseries Modelagem de Data Warehouse}\\[0.6em]
    {\Large\color{graydk} PARTE I --- Modelagem SBD OLTP}

    \vspace{1cm}
    {\Huge\bfseries\color{navy} Relat\'orio}

    \vspace{2cm}

    \begin{tabular}{rl}
      \textbf{Disciplina:} & MAE016 TED-B \textbar\ Big Data \& Data Warehouse \\[0.5em]
      \textbf{Professor:}  & Milton Ramos Ramirez \\[1.5em]
      \textbf{Integrante / DRE:} & Izabela Lima da Silva \quad 124156557 \\[0.3em]
                                 & Caio Meirelles \quad\quad\quad\; 122071557 \\
    \end{tabular}

    \vfill

    {\large 19 de abril de 2026}\\[0.4em]
    {\large UFRJ 2026.1}

    \vspace{0.6cm}
    {\color{navy}\rule{\linewidth}{2pt}}
  \end{center}
\end{titlepage}

% ════════════════════════════════════════════════════════════════════════════
%  SUMÁRIO
% ════════════════════════════════════════════════════════════════════════════
\newpage
\tableofcontents
\newpage

% ════════════════════════════════════════════════════════════════════════════
%  1. INTRODUÇÃO
% ════════════════════════════════════════════════════════════════════════════
\section{Introdu\c{c}\~ao}

Este trabalho tem como objetivo a modelagem de um sistema de banco de dados transacional (OLTP) para uma empresa fict\'icia de loca\c{c}\~ao de ve\'iculos, com p\'atio localizado no Aeroporto Gale\~ao, abrangendo o controle de clientes, frota, reservas, loca\c{c}\~oes, pagamentos e movimenta\c{c}\~ao entre os p\'atios compartilhados de seis empresas independentes de loca\c{c}\~ao de autom\'oveis.

% ════════════════════════════════════════════════════════════════════════════
%  2. UNIVERSO DE DISCURSO
% ════════════════════════════════════════════════════════════════════════════
\section{Universo de Discurso}

O universo de discurso deste sistema compreende as opera\c{c}\~oes de uma locadora de ve\'iculos, incluindo cadastro de clientes, controle de ve\'iculos, gest\~ao de p\'atios, reservas, loca\c{c}\~oes, pagamentos e movimenta\c{c}\~ao de ve\'iculos entre diferentes unidades.

% ════════════════════════════════════════════════════════════════════════════
%  3. REQUISITOS & RELATÓRIOS GERENCIAIS
% ════════════════════════════════════════════════════════════════════════════
\section{Requisitos \& Relat\'orios Gerenciais}

O sistema deve permitir a gera\c{c}\~ao dos seguintes relat\'orios gerenciais:

\begin{itemize}
  \item Quantidade de ve\'iculos por p\'atio
  \item Taxa de ocupa\c{c}\~ao da frota
  \item Hist\'orico de loca\c{c}\~oes por per\'iodo
  \item Movimenta\c{c}\~ao de ve\'iculos entre p\'atios
  \item Disponibilidade de ve\'iculos por categoria
\end{itemize}

% ════════════════════════════════════════════════════════════════════════════
%  4. MODELO CONCEITUAL (MER)
% ════════════════════════════════════════════════════════════════════════════
\section{Modelo Conceitual (MER)}

\begin{figure}[H]
  \centering
  \includegraphics[width=\textwidth, keepaspectratio]{modelo_conceitual.png}
  \caption{Modelo Conceitual}
  \label{fig:conceitual}
\end{figure}

O modelo conceitual foi desenvolvido utilizando o modelo entidade-relacionamento (MER), representando as principais entidades do sistema e seus relacionamentos, como Cliente, Ve\'iculo, Reserva, Loca\c{c}\~ao, P\'atio, Movimenta\c{c}\~ao e Pagamento.

Optou-se por separar as entidades Reserva e Loca\c{c}\~ao, permitindo diferenciar inten\c{c}\~oes de aluguel de sua efetiva\c{c}\~ao. Al\'em disso, a entidade Movimenta\c{c}\~ao foi inclu\'ida para registrar o hist\'orico de deslocamento da frota entre p\'atios.

Principais entidades do sistema:

\begin{itemize}
  \item \textbf{Cliente:} armazena dados dos usu\'arios da locadora
  \item \textbf{Ve\'iculo:} representa os carros dispon\'iveis para loca\c{c}\~ao
  \item \textbf{P\'atio:} locais onde os ve\'iculos est\~ao armazenados
  \item \textbf{Reserva:} solicita\c{c}\~oes pr\'evias de loca\c{c}\~ao
  \item \textbf{Loca\c{c}\~ao:} efetiva\c{c}\~ao do aluguel de um ve\'iculo
  \item \textbf{Pagamento:} registros financeiros das loca\c{c}\~oes
  \item \textbf{Movimenta\c{c}\~ao:} controle de deslocamento de ve\'iculos entre p\'atios
\end{itemize}

% ════════════════════════════════════════════════════════════════════════════
%  5. MODELO LÓGICO
% ════════════════════════════════════════════════════════════════════════════
\section{Modelo L\'ogico}

\begin{figure}[H]
  \centering
  \includegraphics[width=\textwidth, keepaspectratio]{modelo_logico.png}
  \caption{Modelo L\'ogico}
  \label{fig:logico}
\end{figure}

O modelo l\'ogico foi obtido a partir da transforma\c{c}\~ao do modelo conceitual em tabelas relacionais, com defini\c{c}\~ao de chaves prim\'arias e estrangeiras para garantir a integridade referencial.

% ════════════════════════════════════════════════════════════════════════════
%  6. MODELO FÍSICO
% ════════════════════════════════════════════════════════════════════════════
\section{Modelo F\'isico}

O modelo f\'isico foi implementado em SQL seguindo o padr\~ao ANSI SQL (SQL99), com defini\c{c}\~ao das estruturas das tabelas e restri\c{c}\~oes de integridade.

\begin{lstlisting}[style=sql, caption={C\'odigo SQL --- Modelo F\'isico}]
CREATE TABLE cliente (
    id_cliente INT PRIMARY KEY,
    nome VARCHAR(100),
    data_nascimento DATE,
    cpf_cnpj VARCHAR(20),
    cnh VARCHAR(20),
    validade_cnh DATE
);

CREATE TABLE patio (
    id_patio INT PRIMARY KEY,
    endereco VARCHAR(100),
    n_vagas INT
);

CREATE TABLE veiculo (
    id_veiculo INT PRIMARY KEY,
    id_patio INT,
    id_vaga VARCHAR(50),
    modelo VARCHAR(50),
    placa VARCHAR(20),
    FOREIGN KEY (id_patio) REFERENCES patio(id_patio)
);

CREATE TABLE movimentacao (
    id_movimentacao INT PRIMARY KEY,
    id_veiculo INT,
    id_patio_origem INT,
    id_patio_destino INT,
    data_movimentacao DATE,
    FOREIGN KEY (id_veiculo) REFERENCES veiculo(id_veiculo),
    FOREIGN KEY (id_patio_origem) REFERENCES patio(id_patio),
    FOREIGN KEY (id_patio_destino) REFERENCES patio(id_patio)
);

CREATE TABLE reserva (
    id_reserva INT PRIMARY KEY,
    id_cliente INT,
    id_patio_retirada INT,
    categoria_veiculo VARCHAR(50),
    data_inicio DATE,
    data_fim DATE,
    FOREIGN KEY (id_cliente) REFERENCES cliente(id_cliente),
    FOREIGN KEY (id_patio_retirada) REFERENCES patio(id_patio)
);

CREATE TABLE locacao (
    id_locacao INT PRIMARY KEY,
    id_cliente INT,
    id_veiculo INT,
    id_patio_retirada INT,
    id_patio_devolucao INT,
    data_retirada DATE,
    data_devolucao DATE,
    FOREIGN KEY (id_cliente) REFERENCES cliente(id_cliente),
    FOREIGN KEY (id_veiculo) REFERENCES veiculo(id_veiculo),
    FOREIGN KEY (id_patio_retirada) REFERENCES patio(id_patio),
    FOREIGN KEY (id_patio_devolucao) REFERENCES patio(id_patio)
);

CREATE TABLE pagamento (
    id_pagamento INT PRIMARY KEY,
    id_locacao INT,
    data_pagamento DATE,
    valor DECIMAL(10,2),
    FOREIGN KEY (id_locacao) REFERENCES locacao(id_locacao)
);
\end{lstlisting}

Link para reposit\'orio no GitHub:
\url{https://github.com/idevlimes/locadora-oltp}

% ════════════════════════════════════════════════════════════════════════════
%  7. CONSIDERAÇÕES FINAIS
% ════════════════════════════════════════════════════════════════════════════
\section{Considera\c{c}\~{o}es Finais}

O sistema modelado atende aos requisitos propostos, permitindo o controle das opera\c{c}\~oes de loca\c{c}\~ao de ve\'iculos e fornecendo suporte \`a an\'alise gerencial e futura integra\c{c}\~ao com sistemas anal\'iticos.

\end{document}
